\documentclass[12pt,a4paper]{article}

% =========================
%      编码与中文
% =========================
\usepackage[UTF8]{ctex}
\usepackage[T1]{fontenc}
\usepackage{lmodern}

% =========================
%      数学宏包
% =========================
\usepackage{amsmath,amssymb,amsfonts,amsthm,mathtools}
\usepackage{bm}

% =========================
%      版式与超链接
% =========================
\usepackage[a4paper,margin=2.3cm]{geometry}
\usepackage{microtype}
\usepackage{hyperref}
\usepackage{enumitem}

% =========================
%      定理环境
% =========================
\theoremstyle{definition}
\newtheorem{definition}{定义}[section]
\newtheorem{example}{例}[section]
\newtheorem{exercise}{练习}[section]

\theoremstyle{plain}
\newtheorem{theorem}{定理}[section]
\newtheorem{proposition}{命题}[section]
\newtheorem{lemma}{引理}[section]
\newtheorem{corollary}{推论}[section]

% =========================
%      常用符号宏
% =========================
\newcommand{\R}{\mathbb{R}}
\newcommand{\X}{\mathfrak{X}}
\newcommand{\Cinf}{C^{\infty}}
\newcommand{\Id}{\mathrm{Id}}
\newcommand{\dF}{\,\mathrm{d}}
\newcommand{\T}{\mathrm{T}}
\newcommand{\Tp}[1]{\mathrm{T}_{#1}}
\newcommand{\Tps}[1]{\mathrm{T}^{*}_{#1}}
\newcommand{\Tstar}{\mathrm{T}^{*}}
\newcommand{\pair}[2]{\langle #1,#2\rangle}
\newcommand{\GL}{\mathrm{GL}}

\title{余切向量丛（Cotangent Bundle}
\author{}
\date{}

\begin{document}
\maketitle
\tableofcontents

\section{余向量与对偶空间}

\begin{definition}[余向量与对偶空间]
设 $V$ 为有限维实向量空间。\emph{余向量（covector）}是线性泛函 $\omega:V\to\R$。
全体余向量构成向量空间 $V^{*}$，称为 $V$ 的\emph{对偶空间（dual space）}。
\end{definition}

\begin{proposition}[命题 11.1：对偶基]
给定 $V$ 的一组基 $(E_1,\dots,E_n)$，定义 $\varepsilon^1,\dots,\varepsilon^n\in V^*$ 满足
\[
\varepsilon^i(E_j)=\delta^i_j.
\]
则 $(\varepsilon^1,\dots,\varepsilon^n)$ 是 $V^*$ 的一组基，称为 $(E_j)$ 的\emph{对偶基}，并且 $\dim V^*=\dim V$。
\end{proposition}

\begin{exercise}[练习 11.2]
证明命题 11.1。
\end{exercise}

\begin{definition}[对偶映射/转置（dual map / transpose）]
设 $A:V\to W$ 为线性映射。定义 $A^*:W^*\to V^*$：
\[
(A^*\omega)(v)=\omega(Av),\qquad \omega\in W^*,\ v\in V.
\]
\end{definition}

\begin{exercise}[练习 11.3]
证明 $A^*\omega$ 确为 $V$ 上线性泛函且 $A^*$ 是线性映射。
\end{exercise}

\begin{proposition}[命题 11.4：对偶映射性质]
\[
(A\circ B)^* = B^*\circ A^*,\qquad (\Id_V)^*=\Id_{V^*}.
\]
\end{proposition}

\begin{exercise}[练习 11.5]
证明命题 11.4。
\end{exercise}

\begin{corollary}[推论 11.6：反变函子]
把 $V\mapsto V^*$、$A\mapsto A^*$ 的对应给出从实向量空间范畴到自身的\emph{反变函子}。
\end{corollary}

\subsection{第二对偶空间与自然同构}

\begin{definition}[第二对偶空间与自然映射]
定义 $V^{**}=(V^*)^*$。对每个向量空间 $V$，定义自然映射 $\xi:V\to V^{**}$：
\[
\xi(v)(\omega)=\omega(v),\qquad v\in V,\ \omega\in V^*.
\]
\end{definition}

\begin{exercise}[练习 11.7]
(a) 证明 $\xi(v)\in V^{**}$。\quad (b) 证明 $\xi:V\to V^{**}$ 线性。
\end{exercise}

\begin{proposition}[命题 11.8：有限维时 $V\cong V^{**}$]
若 $V$ 有限维，则 $\xi:V\to V^{**}$ 是同构。
\end{proposition}

\begin{definition}[自然配对/尖括号记号]
对 $\omega\in V^*$、$v\in V$，常用 $\pair{\omega}{v}$ 或 $\pair{v}{\omega}$ 表示 $\omega(v)$，称为\emph{自然配对}。
\end{definition}

\section{流形上的切余向量与余切丛}

\begin{definition}[余切空间与切余向量]
设 $M$ 为光滑流形（可带边界），$p\in M$。定义
\[
\Tps{p}M=(\Tp{p}M)^*.
\]
$\Tps{p}M$ 的元素称为 $p$ 处的\emph{切余向量（tangent covector）}。
\end{definition}

\begin{definition}[坐标余标架（coordinate coframe）与分量]
给定坐标图 $(U,(x^i))$，坐标向量基 $\big(\frac{\partial}{\partial x^i}\big|_p\big)$ 在 $\Tps{p}M$ 诱导对偶基，
记为 $(\lambda^i|_p)$（后用 $dx^i|_p$ 表示）。任意 $\omega\in\Tps{p}M$ 可唯一写成
\[
\omega=\omega_i\,\lambda^i|_p,\qquad \omega_i=\omega\Big(\frac{\partial}{\partial x^i}\Big|_p\Big).
\]
\end{definition}

\begin{definition}[余切丛（cotangent bundle）]
定义
\[
\Tstar M=\bigsqcup_{p\in M}\Tps{p}M,
\]
并有自然投影 $\pi:\Tstar M\to M$，$\pi(\omega)=p$（当 $\omega\in\Tps{p}M$）。
\end{definition}

\begin{proposition}[命题 11.9：余切丛是秩 $n$ 的光滑向量丛]
若 $M$ 是 $n$ 维光滑流形，则 $\Tstar M$ 具有唯一的拓扑与光滑结构，使其成为 $M$ 上秩为 $n$ 的光滑向量丛，
且所有坐标余向量场都是光滑局部截面。
\end{proposition}

\begin{exercise}[练习 11.10：一般向量丛的对偶丛]
设 $E\to M$ 为光滑向量丛。定义对偶丛 $E^*\to M$ 的纤维为 $(E_p)^*$。
证明 $E^*\to M$ 为光滑向量丛，且其过渡函数为 $\tau^*(p)=(\tau(p)^{-1})^T$。
\end{exercise}

\begin{definition}[余切丛的自然坐标]
若 $(U,(x^i))$ 为 $M$ 上坐标，则对 $\pi^{-1}(U)$ 上点 $\xi_i\,\lambda^i|_p$，定义坐标
\[
(p,\xi)\ \longmapsto\ \big(x^1(p),\dots,x^n(p),\xi_1,\dots,\xi_n\big)\in \R^{2n}.
\]
称 $(x^i,\xi_i)$ 为 $\Tstar M$ 的\emph{自然坐标}。
\end{definition}

\section{余向量场、余标架与光滑性判别}

\begin{definition}[余向量场/1-形式]
$\Tstar M$ 的截面称为\emph{余向量场}或（微分）\emph{1-形式}。
若不说明，余向量场默认仅连续；光滑者称\emph{光滑余向量场}。
\end{definition}

\begin{definition}[分量函数与收缩（contraction）]
在坐标图中，余向量场可写 $\omega=\omega_i\,\lambda^i$，其中 $\omega_i:U\to\R$ 称为分量函数，
且
\[
\omega_i(p)=\omega_p\Big(\frac{\partial}{\partial x^i}\Big|_p\Big).
\]
若 $X$ 为向量场，则定义函数 $\omega(X):M\to\R$：
\[
\omega(X)(p)=\omega_p(X_p),\qquad \text{坐标下 } \omega(X)=\omega_iX^i.
\]
\end{definition}

\begin{proposition}[命题 11.11：余向量场光滑性判别]
设 $\omega:M\to \Tstar M$ 为粗糙余向量场，则以下等价：\\
(a) $\omega$ 光滑；\\
(b) 在每张坐标图中分量函数 $\omega_i$ 光滑；\\
(c) 每点附近存在坐标图使分量函数光滑；\\
(d) 对每个光滑向量场 $X$，函数 $\omega(X)$ 光滑；\\
(e) 对任意开集 $U$ 上任意光滑 $X$，$\omega(X):U\to\R$ 光滑。
\end{proposition}

\begin{exercise}[练习 11.12]
证明命题 11.11。
\end{exercise}

\begin{definition}[余标架（coframe）]
开集 $U\subseteq M$ 上的\emph{局部余标架}是 $n$ 个余向量场 $(\varepsilon^1,\dots,\varepsilon^n)$，
使得每点 $p\in U$ 的 $\big(\varepsilon^i|_p\big)$ 构成 $\Tps{p}M$ 的一组基。
若 $U=M$，称为整体余标架。
\end{definition}

\begin{example}[例 11.13：坐标余标架]
在坐标图 $(U,(x^i))$ 上的 $(\lambda^i)$（最终记作 $dx^i$）构成局部余标架，称为\emph{坐标余标架}。
\end{example}

\begin{lemma}[引理 11.14：标架与对偶余标架的光滑性等价]
若 $(E_i)$ 是粗糙局部标架，$(\varepsilon^i)$ 为其对偶余标架（$\varepsilon^i(E_j)=\delta^i_j$），
则 $(E_i)$ 光滑当且仅当 $(\varepsilon^i)$ 光滑。
\end{lemma}

\begin{proposition}[命题 11.15：余标架判别光滑性]
设 $(\varepsilon^i)$ 是 $U$ 上光滑余标架，$\omega$ 为粗糙余向量场。
则 $\omega$ 在 $U$ 上光滑 $\Leftrightarrow$ $\omega=\omega_i\varepsilon^i$ 的分量函数 $\omega_i$ 全光滑。
\end{proposition}

\begin{exercise}[练习 11.16]
证明命题 11.15。
\end{exercise}

\begin{definition}[光滑余向量场空间与模结构]
记 $M$ 上全体光滑余向量场为 $\X^*(M)$。对 $f\in \Cinf(M)$、$\omega\in\X^*(M)$，定义
\[
(f\omega)_p=f(p)\omega_p.
\]
则 $\X^*(M)$ 是 $\Cinf(M)$ 上的模。
\end{definition}

\begin{definition}[几何直观：核超平面]
非零余向量 $\omega_p\in \Tps{p}M$ 由其核 $\ker(\omega_p)\subset \Tp{p}M$（余维 1 线性超平面）
以及仿射超平面 $\{v:\omega_p(v)=1\}$（与核平行）刻画，可据此可视化余向量场。
\end{definition}

\section{函数微分}

\begin{definition}[函数的微分]
设 $f\in \Cinf(M)$。定义 1-形式 $df$：
\[
(df)_p(v)=vf,\qquad v\in \Tp{p}M.
\]
\end{definition}

\begin{proposition}[命题 11.18：$df$ 是光滑 1-形式]
若 $f$ 光滑，则 $df$ 是光滑余向量场（光滑 1-形式）。
\end{proposition}

\begin{definition}[坐标表示]
在坐标图中
\[
(df)_p=\frac{\partial f}{\partial x^i}(p)\,dx^i|_p,\qquad
df=\frac{\partial f}{\partial x^i}\,dx^i.
\]
并且 $dx^j=\lambda^j$（坐标余向量场即坐标函数的微分）。
\end{definition}

\begin{example}[例 11.19]
若 $f(x,y)=x^2y\cos x$，则
\[
df=(2xy\cos x-x^2y\sin x)\,dx + x^2\cos x\,dy.
\]
\end{example}

\begin{proposition}[命题 11.20：微分的性质]
设 $f,g\in\Cinf(M)$：\\
(a) $d(af+bg)=a\,df+b\,dg$（$a,b$ 常数）；\\
(b) $d(fg)=f\,dg+g\,df$；\\
(c) 在 $g\neq 0$ 上，$d(f/g)=\frac{g\,df-f\,dg}{g^2}$；\\
(d) 若 $h:J\to\R$ 光滑且 $f(M)\subset J$，则 $d(h\circ f)=(h'\circ f)\,df$；\\
(e) $f$ 常数 $\Rightarrow df=0$。
\end{proposition}

\begin{exercise}[练习 11.21]
证明命题 11.20。
\end{exercise}

\begin{proposition}[命题 11.22：微分为零的函数]
若 $f\in\Cinf(M)$，则 $df=0$ 当且仅当 $f$ 在 $M$ 的每个连通分支上为常数。
\end{proposition}

\begin{theorem}[命题 11.23：沿曲线的导数公式]
设 $\gamma:J\to M$ 光滑曲线、$f\in\Cinf(M)$，则
\[
(f\circ\gamma)'(t)=(df)_{\gamma(t)}\big(\gamma'(t)\big).
\]
\end{theorem}

\begin{exercise}[练习 11.24]
证明：$p$ 是 $f$ 的临界点当且仅当 $(df)_p=0$。
\end{exercise}

\section{余向量场的拉回与限制}

\subsection{逐点拉回与余向量场拉回}

\begin{definition}[逐点拉回/余切映射]
设 $F:M\to N$ 光滑。对 $p\in M$，微分
\[
dF_p:\Tp{p}M\to \Tp{F(p)}N
\]
诱导对偶映射（逐点拉回）
\[
dF_p^*:\Tps{F(p)}N\to \Tps{p}M,\qquad dF_p^*(\omega)(v)=\omega(dF_p(v)).
\]
\end{definition}

\begin{definition}[余向量场拉回]
若 $\omega$ 为 $N$ 上余向量场，定义 $M$ 上余向量场 $F^*\omega$：
\[
(F^*\omega)_p = dF_p^*\big(\omega_{F(p)}\big),\qquad
(F^*\omega)_p(v)=\omega_{F(p)}(dF_p(v)).
\]
\end{definition}

\begin{proposition}[命题 11.25：拉回的两条基本性质]
设 $u$ 为 $N$ 上连续函数、$\omega$ 为 $N$ 上余向量场，则
\[
F^*(u\omega)=(u\circ F)\,F^*\omega.
\]
若进一步 $u$ 光滑，则
\[
F^*du = d(u\circ F).
\]
\end{proposition}

\begin{proposition}[命题 11.26：拉回保持连续/光滑]
若 $\omega$ 连续，则 $F^*\omega$ 连续；若 $\omega$ 光滑，则 $F^*\omega$ 光滑。
并且在坐标中可写为
\[
F^*\omega=(\omega_j\circ F)\,d(y^j\circ F)=(\omega_j\circ F)\,dF^j.
\]
\end{proposition}

\begin{example}[例 11.27]
令 $F:\R^3\to\R^2$，$F(x,y,z)=(u,v)=(x^2y,\ y\sin z)$，并令
\[
\omega=u\,dv+v\,du.
\]
则
\[
F^*\omega
=2xy^2\sin z\,dx+2x^2y\sin z\,dy+x^2y^2\cos z\,dz.
\]
\end{example}

\begin{example}[例 11.28：极坐标替换法]
在右半平面 $H=\{(x,y):x>0\}$ 上用极坐标 $(r,\theta)$：$x=r\cos\theta,\ y=r\sin\theta$。
则
\[
x\,dy-y\,dx = r^2\,d\theta.
\]
\end{example}

\subsection{限制到子流形}

\begin{definition}[限制（restriction）为拉回]
设 $i:S\hookrightarrow M$ 为浸入子流形的包含映射。对 $M$ 上余向量场 $\omega$，其在 $S$ 上的限制定义为
\[
i^*\omega,
\]
其作用满足 $(i^*\omega)_p(v)=\omega_p(v)$（$v\in\Tp{p}S$ 视为 $\Tp{p}M$ 的子空间向量）。
\end{definition}

\begin{example}[例 11.29：限制可能为零]
在 $\R^2$ 上取 $\omega=dy$，令 $S$ 为 $x$ 轴（嵌入子流形），则 $\omega$ 在 $\R^2$ 处处非零，
但 $i^*\omega\equiv 0$，因为 $i^*dy=d(y\circ i)=0$。
\end{example}

\begin{exercise}[练习 11.30]
证明 $d(f|_S)=i^*(df)$。推出：$df$ 拉回到 $S$ 为零当且仅当 $f$ 在 $S$ 每个连通分支上为常数。
\end{exercise}

\section{线积分}

\subsection{区间上的积分与换元不变性}

\begin{definition}[区间上的 1-形式积分]
在紧区间 $[a,b]\subset\R$ 上，若 $\omega=f(t)\,dt$，定义
\[
\int_{[a,b]}\omega := \int_a^b f(t)\,dt.
\]
\end{definition}

\begin{proposition}[命题 11.31：递增微分同胚不变性]
若 $\varphi:[c,d]\to[a,b]$ 为递增微分同胚，则
\[
\int_{[c,d]}\varphi^*\omega = \int_{[a,b]}\omega.
\]
\end{proposition}

\begin{exercise}[练习 11.32]
若 $\varphi$ 递减，则 $\int_{[c,d]}\varphi^*\omega = -\int_{[a,b]}\omega$。
\end{exercise}

\subsection{曲线段、线积分与性质}

\begin{definition}[曲线段/分段光滑曲线段]
曲线段指连续曲线 $\gamma:[a,b]\to M$（$[a,b]$ 紧）。若 $\gamma$ 光滑（将区间视为带边界流形），称\emph{光滑曲线段}。
若存在有限分割 $a=a_0<\cdots<a_k=b$ 使得 $\gamma|_{[a_{i-1},a_i]}$ 光滑，则称\emph{分段光滑曲线段}。
\end{definition}

\begin{proposition}[命题 11.33：连通流形上两点可由分段光滑曲线连接]
若 $M$ 连通，则任意两点可由一条分段光滑曲线段连接。
\end{proposition}

\begin{definition}[线积分]
若 $\gamma:[a,b]\to M$ 光滑且 $\omega\in\X^*(M)$，定义
\[
\int_{\gamma}\omega := \int_{[a,b]}\gamma^*\omega.
\]
若 $\gamma$ 分段光滑，则在各光滑子区间上求和定义 $\int_\gamma\omega$。
\end{definition}

\begin{proposition}[命题 11.34：线积分性质]
设 $\gamma$ 分段光滑、$\omega,\omega_1,\omega_2\in\X^*(M)$：\\
(a) 线性：$\int_\gamma(c_1\omega_1+c_2\omega_2)=c_1\int_\gamma\omega_1+c_2\int_\gamma\omega_2$；\\
(b) 常曲线：$\gamma$ 常值 $\Rightarrow \int_\gamma\omega=0$；\\
(c) 拼接可加：若 $\gamma=\gamma_1*\gamma_2$，则 $\int_\gamma\omega=\int_{\gamma_1}\omega+\int_{\gamma_2}\omega$；\\
(d) 与拉回相容：$\int_\gamma F^*\eta=\int_{F\circ\gamma}\eta$。
\end{proposition}

\begin{exercise}[练习 11.35]
证明命题 11.34。
\end{exercise}

\begin{example}[例 11.36：绕圈积分]
在 $M=\R^2\setminus\{0\}$ 上取
\[
\omega=\frac{x\,dy-y\,dx}{x^2+y^2},
\quad \gamma(t)=(\cos t,\sin t),\ t\in[0,2\pi],
\]
则 $\displaystyle \int_\gamma \omega=2\pi$。
\end{example}

\begin{definition}[重参数化]
若 $\tilde\gamma=\gamma\circ\varphi$，其中 $\varphi:[c,d]\to[a,b]$ 为微分同胚，则称 $\tilde\gamma$ 为 $\gamma$ 的重参数化；
$\varphi$ 递增/递减分别称正向/反向重参数化。
\end{definition}

\begin{proposition}[命题 11.37：线积分与参数无关]
若 $\tilde\gamma$ 为 $\gamma$ 的重参数化，则
\[
\int_{\tilde\gamma}\omega=
\begin{cases}
\int_\gamma\omega,& \tilde\gamma\text{ 正向重参数化},\\
-\int_\gamma\omega,& \tilde\gamma\text{ 反向重参数化}.
\end{cases}
\]
\end{proposition}

\begin{proposition}[命题 11.38：普通积分表达式]
若 $\gamma:[a,b]\to M$ 分段光滑，则
\[
\int_\gamma\omega=\int_a^b \omega_{\gamma(t)}\big(\gamma'(t)\big)\,dt.
\]
\end{proposition}

\begin{theorem}[定理 11.39：线积分基本定理]
设 $f\in\Cinf(M)$，$\gamma:[a,b]\to M$ 分段光滑，则
\[
\int_\gamma df = f(\gamma(b)) - f(\gamma(a)).
\]
\end{theorem}

\section{恰当/保守/闭 1-形式与庞加莱引理}

\begin{definition}[恰当（exact）与势函数]
若 $\omega\in\X^*(M)$，存在 $f\in\Cinf(M)$ 使 $\omega=df$，则称 $\omega$ \emph{恰当}，$f$ 称为 $\omega$ 的\emph{势函数}。
\end{definition}

\begin{definition}[保守（conservative）]
若对任意分段光滑闭曲线段 $\gamma$ 都有 $\int_\gamma\omega=0$，则称 $\omega$ \emph{保守}。
\end{definition}

\begin{proposition}[命题 11.40：路径无关刻画保守性]
$\omega$ 保守 $\Longleftrightarrow$ 任意起终点相同的分段光滑曲线段 $\gamma,\tilde\gamma$ 满足
\[
\int_\gamma\omega=\int_{\tilde\gamma}\omega.
\]
\end{proposition}

\begin{exercise}[练习 11.41]
证明命题 11.40。
\end{exercise}

\begin{theorem}[定理 11.42：保守当且仅当恰当]
在带边界或不带边界的光滑流形 $M$ 上，光滑余向量场 $\omega$ 保守 $\Longleftrightarrow$ 恰当。
\end{theorem}

\begin{example}[例 11.43：穿孔平面上闭但不恰当]
例 11.36 的 $\omega=\frac{x\,dy-y\,dx}{x^2+y^2}$ 在 $\R^2\setminus\{0\}$ 上不是恰当的，
因为存在闭曲线使 $\int_\gamma\omega=2\pi\neq 0$，故不保守。
\end{example}

\subsection{闭（closed）与等价刻画}

\begin{definition}[闭 1-形式（closed）]
在坐标中写 $\omega=\omega_i\,dx^i$。若在每张坐标图中对所有 $i,j$ 成立
\[
\frac{\partial\omega_j}{\partial x^i}=\frac{\partial\omega_i}{\partial x^j},
\]
则称 $\omega$ \emph{闭}。
\end{definition}

\begin{proposition}[命题 11.44：恰当必闭]
每个恰当余向量场都是闭的。
\end{proposition}

\begin{proposition}[命题 11.45：闭性的坐标无关刻画]
对光滑余向量场 $\omega$，以下等价：\\
(a) $\omega$ 闭；\\
(b) 每点附近存在坐标图使其分量满足上述偏导对称条件；\\
(c) 对任意开集 $U\subseteq M$ 及任意 $X,Y\in\X(U)$，
\[
X(\omega(Y)) - Y(\omega(X)) = \omega([X,Y]).
\]
\end{proposition}

\begin{corollary}[推论 11.46：局部微分同胚下闭/恰当的保持性]
若 $F:M\to N$ 为局部微分同胚，则拉回 $F^*:\X^*(N)\to\X^*(M)$
把闭 1-形式拉回为闭 1-形式，把恰当 1-形式拉回为恰当 1-形式。
\end{corollary}

\begin{example}[例 11.47：闭与非闭的对比]
在 $\R^2$ 上
\[
\omega=y\cos(xy)\,dx+x\cos(xy)\,dy
\]
满足闭性条件（并可猜测 $\omega=d(\sin(xy))$）。
而
\[
\eta=x\cos(xy)\,dx+y\cos(xy)\,dy
\]
不闭，因此不恰当。
\end{example}

\begin{example}[例 11.48：域的形状影响“闭是否恰当”】【全局性]
例 11.36 的 $\omega$ 在 $\R^2\setminus\{0\}$ 上闭但不恰当；
若限制到右半平面 $U=\{(x,y):x>0\}$，则可写成
\[
\omega=d(\tan^{-1}(y/x))\quad(\text{极坐标中 } \omega=d\theta),
\]
显示“闭 $\Rightarrow$ 恰当”依赖定义域拓扑（是否有洞）。
\end{example}

\subsection{星形域与庞加莱引理}

\begin{definition}[星形开集（star-shaped）]
设 $U\subset V$（$V$ 为有限维向量空间）。若存在 $c\in U$ 使得对任意 $x\in U$，线段 $[c,x]\subset U$，
则称 $U$ 关于 $c$ 星形。每个凸集都是星形的。
\end{definition}

\begin{theorem}[定理 11.49：余向量场的庞加莱引理]
若 $U$ 是 $\R^n$ 或 $\mathbb{H}^n$ 的星形开子集，则 $U$ 上每个闭余向量场都是恰当的。
\end{theorem}

\begin{corollary}[推论 11.50：闭 1-形式局部恰当]
设 $\omega$ 是 $M$ 上闭余向量场，则 $M$ 的每个点都有邻域 $U$ 使得 $\omega|_U$ 恰当（存在 $f\in\Cinf(U)$ 使 $\omega|_U=df$）。
\end{corollary}

\subsection{求势函数的计算法示例}

\begin{example}[例 11.51：计算势函数]
在 $\R^3$ 上令
\[
\omega=e^{y^2}\,dx+2xye^{y^2}\,dy-2z\,dz.
\]
其势函数 $f$ 需满足
\[
\frac{\partial f}{\partial x}=e^{y^2},\quad
\frac{\partial f}{\partial y}=2xye^{y^2},\quad
\frac{\partial f}{\partial z}=-2z.
\]
可取一势函数
\[
f(x,y,z)=x e^{y^2}-z^2,
\]
其余势函数与之相差常数。
\end{example}

\end{document}